\documentclass[10pt,a4paper]{article}

% --- XeLaTeX: fonts + Russian ---
\usepackage{fontspec}
\usepackage{polyglossia}
\setdefaultlanguage{russian}
\setotherlanguage{english}
\setmainfont{PT Sans}
\setmonofont[Scale=MatchLowercase]{PT Mono}
\newfontfamily\cyrillicfonttt[Scale=MatchLowercase]{PT Mono}

% --- layout ---
\usepackage[a4paper,top=0.7cm,bottom=0.7cm,left=1.3cm,right=1.3cm]{geometry}
\usepackage{titlesec}
\usepackage{enumitem}
\usepackage{xcolor}
\usepackage[hidelinks]{hyperref}

\definecolor{rule}{HTML}{333333}

\titleformat{\section}{\normalsize\bfseries}{}{0em}{\MakeUppercase}[{\color{rule}\titlerule[0.6pt]}]
\titlespacing*{\section}{0pt}{4pt}{2.2pt}

\setlist[itemize]{leftmargin=1.15em, labelsep=0.5em, itemsep=1.2pt, topsep=2pt, parsep=0pt, before=\vspace{1pt}}
\renewcommand{\labelitemi}{\textbullet}

\setlength{\parindent}{0pt}
\setlength{\parskip}{0pt}
\pagestyle{empty}
\linespread{0.96}
\tolerance=1500
\emergencystretch=2.5em


\newcommand{\role}[2]{\textbf{#1}\hfill #2\par}

\begin{document}

% ===== Header =====
\begin{center}
  {\LARGE\bfseries Михайловский Илья Евгеньевич}\\[2pt]
  {\large Data Scientist, NLP / LLM}\\[3pt]
  Москва \textperiodcentered{} офис, гибрид или удалённо \textperiodcentered{} английский C1\\[2pt]
  Telegram: \href{https://t.me/Iluuuaaa}{@Iluuuaaa} \textperiodcentered{}
  \href{mailto:ilyamihailovsy@gmail.com}{ilyamihailovsy@gmail.com} \textperiodcentered{}
  \href{https://github.com/iluuua}{github.com/iluuua}
\end{center}

% ===== Summary =====
\section{О себе}
Учусь на программе <<Прикладная математика и информатика>> в Московском политехе, выпуск в 2029 году.
С января 2026 года разрабатываю NLP/LLM-часть ReachFlow: классификацию диалогов, извлечение фактов,
агентные сценарии, память и векторный поиск. Обучал модели с нуля на PyTorch.
Финалист AIDAO 2025: IoU 0.5606 на скрытом тесте при 0.564 у победителя.

\section{Навыки}
\begin{itemize}
  \item NLP и LLM: классификация текста, извлечение фактов, prompt engineering,
        ответы по JSON Schema, function calling, MCP, память агентов, эмбеддинги,
        pgvector и HNSW, RAG, выбор моделей и резервных провайдеров, оценка стоимости и латентности
  \item ML и DL: PyTorch, NumPy, pandas; обучение с нуля и transfer learning, mixed precision,
        функция потерь под целевую метрику, калибровка порога, дисбаланс классов;
        ResNet, FPN, сегментация, IoU
  \item Инженерия: Python, FastAPI, asyncio, PostgreSQL и SQL, Redis, Docker, CI/CD,
        Linux, Git, pytest
  \item Математика и алгоритмы: теория вероятностей, матстатистика, линейная алгебра,
        методы оптимизации; C++, алгоритмы и структуры данных
\end{itemize}

\section{Опыт работы}
\role{ReachFlow, Founder}{январь 2026 - настоящее время}
\textit{B2B SaaS для работы с обращениями и диалогами в Telegram \textperiodcentered{} \href{https://reachflow.tech}{reachflow.tech}}
Основатель продукта; отвечаю за NLP/LLM и backend.
\begin{itemize}
  \item Разработал классификацию состояния диалога по 16 типам событий и 8 классам настроения.
        Каждое событие связано с исходным сообщением и оценкой уверенности.
        Факты о собеседнике сохраняются по 18 типизированным ключам с версионированной
        политикой памяти и порогами хранения.
  \item Собрал агента вокруг реестра \texttt{ToolSpec} с 43 инструментами.
        Из реестра формируются описания для промпта, проверки прав по уровню риска
        и \texttt{inputSchema} MCP-эндпоинта. Ограничил бюджет вызовов инструментов;
        при эскалации диалог передаётся менеджеру.
  \item Настроил выполнение 25 LLM-задач с 22 схемами Pydantic (\texttt{extra=forbid}).
        Схемы адаптируются под strict-режим провайдеров, ответы проверяются на сервере.
        При нарушении схемы возможна повторная попытка или переход к резервному провайдеру.
  \item Реализовал косинусный поиск на pgvector с HNSW-индексом и эмбеддингами
        \textit{text-embedding-3-small} размерности 1536. При разборе семантического фильтра
        выяснил, что он отклонял 55 кандидатов из 63 из-за формата и длины текста.
        Перевёл фильтр в режим наблюдения.
  \item Объединил пять классификаторных вызовов в два. Объём входного текста на этом этапе
        уменьшился на 14\% для короткого диалога и на 43\% для длинного;
        каждая секция сохранила отдельную обработку ошибок.
        Настроил трассировку модели, времени выполнения, повторных попыток и стоимости.
\end{itemize}
\vspace{2pt}

\role{Delnovo, Chief Technology Officer}{июль 2023 - май 2026}
\begin{itemize}
  \item Разрабатывал коммерческих ботов для поддержки и продаж. Для дистрибьютора пластиковых окон
        собрал RAG на PostgreSQL с векторным поиском по базе знаний и интеграцией с Битрикс24.
        По оценке заказчика бот автоматизировал около 80\% типовых консультаций.
\end{itemize}
\vspace{2pt}

\role{NEAR AI, Solutions Engineer}{июнь 2023 - июль 2024}
\begin{itemize}
  \item Решал алгоритмические задачи на C++ для обучения децентрализованного ИИ:
        оптимизировал время и память, оценивал асимптотику, проверял код участников международной команды.
\end{itemize}

\section{ML-проекты и соревнования}
\role{AIDAO 2025, финалист}{ноябрь 2025}
Олимпиада ФКН НИУ ВШЭ и Яндекс Образования. Топ-30 команд из 248 заявок от участников
из 14 стран; очный финал, 32 часа. Задача команды беспилотных технологий Яндекса:
occupancy-карта в bird's-eye view по камерам автомобиля, метрика IoU.
\begin{itemize}
  \item Реализовал Lift-Splat-Shoot с нуля на PyTorch: общий для камер ResNet-50 с FPN,
        depth-head на 64 интервала глубины с обучаемой температурой softmax,
        дифференцируемая проекция признаков в BEV-сетку 188$\times$126.
  \item Основной прирост дала покадровая калибровка. Функция потерь: BCE с
        \texttt{pos\_weight} по балансу классов и дифференцируемый surrogate mIoU.
        IoU на скрытом тесте: 0.5606; у победителя: 0.564.
        Код: \href{https://github.com/iluuua/aidao-2025-bev-occupancy}{aidao-2025-bev-occupancy}.
\end{itemize}

\role{Лидеры цифровой трансформации 2026, задача №5}{2026}
Капитан команды xixi-square. Разработал алгоритм обнаружения препятствий по облакам точек
3D-лидара Pandar128 и интеграцию с ROS 2 Humble. Геометрия тоннеля и рельсов задаёт габарит
движения; препятствия подтверждаются по нескольким кадрам. Решение упаковано в Docker.
На синтетике организаторов найдены все 8 препятствий; на проверенных 15,7 км реальных записей
ложных тревог не было.

\section{Образование и курсы}
\textbf{Московский политехнический университет}, прикладная математика и информатика, выпуск 2029.
\textbf{НИУ ВШЭ, ФКН}, вольнослушатель курсов той же программы.
\textbf{Deep Learning School} ФПМИ МФТИ, 2025.
\textbf{ОЦ <<Сириус>>}, <<Январская научная школа по математике и программированию>>,
92 академических часа, январь 2025.
\textbf{Т-Образование}, <<Т-Поколение>>, курс <<Алгоритмы>>, параллель B, 2025.

\end{document}
